\lesson{5}{Nov 16 2021 Tue (14:22:49)}{Media and Politics}{Unit 3}

\subsubsection*{Media}

This refers to all the vehicles or methods through which a wide audience receives \bf{Messages} and \bf{Information}. It includes print media such as:

\begin{itemize}
    \item Newspapers
    \item Magazines
    \item Books
\end{itemize}

A wide range of electronic media has emerged including:

\begin{itemize}
    \item Digital music
    \item Music and Movies
    \item Smart Phones
    \item DVDs and CDs
    \item Digital Photographs
    \item Video Games
\end{itemize}

Electronic media that is used to transmit information to the public, like:

\begin{itemize}
    \item Television
    \item Radio
\end{itemize}

is also known as \bf{Broadcast Media}. Certainly various other visual arts can send a message to people, such as:

\begin{itemize}
    \item Paintings
    \item Roadway Billboards
    \item Bumper Stickers
\end{itemize}

The internet may deserve a media category all of its own, since it combines so many different types of media. Any reputable news organization or media outlet would certainly have a presence on the World Wide Web today.

The term \it{"Mass Media"} usually refers to the news organizations and agencies that control images and audio delivered publicly through the various technologies. Newspaper offices, TV stations, and e-book publishers are just a few of these businesses. Media could seek to reach a local, regional, national, or international audience. It can focus on a particular type of message, such as political campaigns, or address a wide range of interests and topics. In the United States, the media have a key role in the relationship between the people and their government.

\subsubsection*{Role of the Media in Politics}

In a \bf{Representative Democracy}, the media have an important role in helping \bf{Elected Officials} and \bf{Constituents} communicate with each other. \bf{Candidates} running for office explain their platform and debate policy with other candidates through the media. A heavy cost of campaigning is the purchase of media time or space for candidates to relay their message to voters. They use multiple broadcast and print organizations to reach as wide a voting audience as possible.

The work of \bf{Congress} and \bf{State Legislatures} is available on television, in print records, and on Internet websites. Voters have the ability to review the work of their elected officials at any time. They can scrutinize the actions and voting record of a public official before choosing whether to support him or her during their re-election campaign. The public can use the media to send messages to government which relate how they feel about certain policies or about the leaders themselves. Communication from \bf{Constituents} influences a representative's focus while in office.

The media are not just a tool for the public or for \bf{Elected Leaders} to develop a relationship with the other. Those who work in the media have their own opinions and provide commentary on government-related activities and the reactions of the public. Choices made in the content and style of a political message can have a powerful effect on voters and officials alike.

In $2003$, top officials in the White House engaged the public in a series of live chats, answering questions about President Bush's administration. Dan Bartlett, White House Director of Communications at the time, participates in one of these chats from his office in the White House. In the years since, presidents have used social media to directly connect with the public.

\subsubsection*{The forming of the Media}

\paragraph{Pre-1800} Several important developments related to \bf{Print Media} occurred prior to the founding of the \bf{United States}. In the $1400s$, \it{Johan Gutenberg} developed movable type, making the publishing of large volumes of newspapers and books possible. The first daily newspaper, the advertisement in print, the \bf{Postal Service}, and a written journal geared for a specific audience appeared in the $1600s$. \it{Benjamin Franklin} is a famous American of the $1700s$ who started the first circulating library and published the first political cartoon. That century also saw the first experiments with wood-based paper and steel point pens to replace feather quills. \bf{Print Media} was the dominant form for political communication in the early \bf{United States}, including transcribed speeches of officials. The \it{First Amendment} was created in part to protect the rights of free speech and press.

\paragraph{1800s} The mid-$1800s$ saw an explosion of new American newspapers, magazines, and novels. Print media was still dominant, and newspaper editors and writers used their media to share the words of important leaders and to send political messages of their own on specific topics.

\paragraph{1890s} Throughout the $1800s$, developments such as photographic film and the telegraph paved the way for radio and film in the $1890s$. The ability to see realistic images and hear the voices of political leaders would help the public relate to them. Both would grow to become primary media for delivery of general news and information about politics through the early 1900s.

\paragraph{1930s-1960s} Television was developed in the $1930s$ and was revealed to an awed world audience at the $1939$ World's Fair. TV ownership among Americans would not be common for several more years. Yet soon many realized its potential not only for entertainment and commercial advertising, but also political advertising. As television ownership in the United States rose, so did television's role in politics and elections.

\paragraph{1980s} In the $1980s$, digitized recordings were developed. Compact discs, or CDs, began to replace cassette tapes. Videocassette recorders, or VCRs, made possible a wide exchange of recorded sound and video. Soon the DVD, or digital video disc, would replace the video cassette. Not only could people purchase their favorite movies and music, they could now record broadcasts they otherwise would miss due to work or family obligations. Political communications such as the annual State of the Union address or political debates could now be viewed at one's leisure. Of course, you had to have the required equipment.

\paragraph{1990s} The internet, developed in the late $1960s$, was made available to the public audience in the early $1990s$ via the internet. These inventions would transform and diversify the nature of political communications in the \bf{United States}. While the internet expanded the increasing use of technology that was already available, the most significant achievement was the development of social networking worldwide. The instant access to information and ideas has led to greater scrutiny of elected officials and public policies. The later years of this decade saw an explosion of cell phone users. Today, most Americans own a smartphone rather than a cell phone. Candidates for office, government agencies, and elected officials continue to explore and leverage new technologies as rapidly as the American public embraces them. Recent presidents have made use of social media to connect directly with the public.

\subsubsection*{State of the Union Address}

The Constitution states that one of the duties of the president is to deliver an annual message to Congress about the state of the union. In 1801, Thomas Jefferson wrote his message to Congress, starting the tradition of delivering the information in writing. President Calvin Coolidge was the first to have his speech broadcast over radio in 1923. The 1947 address by President Harry S. Truman was the first to be televised. President Clinton's speech was the first broadcasted over the World Wide Web in 1997.

While primarily intended for Congress, it soon became an important medium for a president to explain his policy agenda for the nation in the coming year. The speech gives representatives and senators an idea of the policy changes the president may ask for. It also helps inform the public of these changes, bringing attention to a wide variety of domestic and foreign policy issues. Now, during, and immediately after the president delivers the speech each year, Americans can access videos and text transcriptions of the speech over the internet.

\paragraph{Fact Checking the President} Each year the State of the Union Address is scrutinized for "spin." Spin is the phrasing of information and facts in a way that benefits the speaker's interests. Presidents usually announce that the state of the union is "strong." They credit their own efforts for that strength and outline their plans to make it even better. To be fair, the analyses of the president's speech can contain just as much spin themselves, whether from opposing political groups or reputable news media. Below, examine the facts relating to statements of President Obama in his first State of the Union Address.

\subsubsection*{Bias Affecting Political Communication}

The media are, naturally, made up of people who have their own ideas and perspectives about the way politics should work. Those responsible for creating, delivering, or distributing political messages can shape or alter their impact in different ways. Some examples include a campaign ad, televised speech, personal interview, or other various formats. During or after a campaign debate, reporters will discuss on television, radio, and the internet their opinions on what the candidates had to say. The commentary can range from objective analysis to overt bias.

Bias is a preference toward a particular belief or attitude. Often people think of "bias" as a negative thing, especially when in the context of stereotypes of a general group of people. However, people naturally have preferences, and bias can be positive or neutral, instead of negative. For instance, you could have a bias towards a favorite flavor or brand of ice cream. A preference of one brand is not negative but does show a bias for one thing over another. The same is true for biases concerning an individual's political ideas.

Bias in political messages delivered by government or the public through the media can manifest in different ways. It can be obvious. For example, you would expect a Republican senator to support a conservative policy agenda. Yet some messages are much more subtle. You might hear people refer to a particular TV channel as conservative or liberal. People who work for it may not directly say they favor a particular bias or that their bosses want them to favor a particular bias, but we can discover bias by the nature of their various stories and attention to different policies and officials.

\begin{definition}[Bias]
    Bias is a preference toward a particular belief or attitude. Political messages, by their very nature, exhibit bias. Whether a speech, advertisement, or other type, a political message is trying to convince people to support a policy idea, group, or candidate for office. A political advertisement, such as an ad promoting a candidate for office, has bias because the maker shows preference toward a particular belief or attitude. Bias may be evident in a political message through its different forms, such as accuracy, omission, and emotional appeal.
\end{definition}

\begin{definition}[Types of Bias-Accuracy]
    Often, the strongest persuasive arguments use facts and statistics to back up their position. Those that disagree with the position have a harder time arguing against it in the face of research and evidence. However, political messages are sometimes vague about facts or may distort them.
    
    For example, during a debate a candidate may point out that his opponent failed to stop an increase of crime while serving as the local police chief. Perhaps that's not the whole story. The former police chief points out that while the specific number of crimes rose during her tenure, so did the population. She defends his attack by pointing out data that shows the number of crimes per person did in fact decrease. By making a general statement that crime increased under the police chief's watch, the first candidate has distorted the facts in an attempt to sway voters' support away from his opponent.
\end{definition}

\begin{definition}[Types of Bias–Omission]
    Another type of bias that is similar to accuracy is the strategy of omission. Omission happens when a political statement or advertisement leaves out key information related to the contents of the message.
    
    For example, the opponents of legislators running for re-election often point out when the legislator has voted against a popular policy idea, especially when that vote seems out of line with the legislator's political party or ideology. What the opponent's message may have left out are the reasons why the legislator voted against it. Perhaps there were two similar bills on the issue and the legislator thought the other one would be better. Alternatively, perhaps a small part of the proposed bill was problematic, and the legislator hoped to kill that bill and draft an improved one. Whatever the reasons, purposeful omission of details is used in political media in an attempt to influence voters, and candidates sometimes are forced to explain their omissions in debates or media interviews.
\end{definition}

\begin{definition}[Types of Bias–Emotional]
    Appeal–Besides the selective choosing of facts and details, political messages can sway voters by appealing to their emotions. Statements and visual images that appeal to our personal values, beliefs, and desires can elicit powerful emotional reactions, either positive or negative.
    
    For example, political candidates will often pose for photographs with their children or grandchildren. Without stating it obviously, it sends a message to capture the emotional appeal that the candidate values family and children, a common value held by most Americans. Interest groups may run ads describing the consequences they predict will come if a certain policy idea is not changed. For example, an animal rights group might show images from farms or slaughterhouses that people find disturbing to convince them to demand laws that would increase protection for animal rights.
\end{definition}

\begin{definition}[Symbols]
    Symbols are images or text in advertisements that represent something else or have a meaning not obvious from the picture itself. For example, a donkey is the symbol for the Democratic Party while an elephant is the symbol for the Republican Party. The symbols originated in political cartoons of the 1800s. Today, the donkey and elephant are rather benign, though symbols can reflect bias and political messages.
    
    For example, candidate ads that use an American flag are trying to convince viewers that the candidate is proud of his or her country. The flag stands in as a symbol for patriotism. Another example of the symbol in political messages is clothing. For example, when appearing before a group of mineworkers to discuss her plans to increase benefits for them, a candidate is not likely to wear a business suit. She may wear jeans and a shirt with rolled up sleeves, sending a subtle message that she can relate to the mineworkers and their concerns.
\end{definition}

\begin{definition}[Propaganda]
    Propaganda is information an organization or government sends out to promote a policy, idea, or cause. In general, it means any type of political message with a bias. An advertisement is a type of propaganda. Many people think of the word propaganda as something bad, but its original meaning is neutral. Propaganda can promote both popular and unpopular ideas.
\end{definition}

\subsubsection*{Regulation Affecting Political Communications}

Regulatory policies affect the nature of political communications, with the goal of protecting \it{equal access, objectivity, and competition}. Since its founding, the \bf{Federal Communications Commission} has operated as police of the airwaves, enforcing decency standards like limitations on use of profanity during the hours when children are more likely to be in the audience. The executive agency also created rules for political campaigns.

The \bf{Federal Communications Commission} expected the media to perform a public service by being a source of information on critical issues within the community, such as disaster preparedness. Regulations prevented a media organization from buying out the majority of limited radio and broadcast channels and consolidating toward monopoly. This protected viewers' rights of choice in programming, as well as ensuring access to multiple perspectives and news sources. Media ownership regulations were common throughout the twentieth century, such as inhibiting one company from owning both a newspaper and a TV station in the same market at the same time. \bf{Congress} relaxed or eliminated these regulations in the $1980s$ and $1990s$. They still exist, but the \bf{Federal Communications Commission} is not required to enforce them.

The \bf{Telecommunications Act} of $1996$ in particular relaxed ownership rules. Today, $6$ corporations own the majority of \it{TV and radio airwaves, movie production companies, publishing companies, and well-known internet sites}. The \bf{Walt Disney Company} is one of the corporations, owning companies such as \bf{ESPN television channel} and \bf{Pixar Animation Studios}.

\newpage
